\section{Eleven integers, degree 352: saturating by a hidden square root}
\label{sec:integer352}
The primary equation below has eleven unrestricted integer unknowns, exact
total degree $352$, and $132$ live binary arithmetic operations. It improves
the incoming degree-$416$ endpoint in degree, at a cost of seven operations.
Three new nondominated companions have degree/cost pairs
$428/123$, $376/127$, and $364/131$. The six fixed source34 numerals are
unchanged. This is a new ordinary arithmetic reduction with the inherited
native/compiler/refutation calibration, not a consistency proof, a global
minimum, or a new end-to-end proof-assistant certificate.

\subsection{The complete polynomial}
The unknowns are
\[
 (j,C,z,h,r,v,w,T,c,n,s)\in\mathbb Z^{11}.
\]
Use the fully authenticated numerals
$\beta,K_0,M_C,M_F,A_c,B_c$, with
$\beta=2^{1064921}-1$, $A_c=2^{1113}$, $B_c=2^{666104}$ and the already
shifted $M_F$. Every symbol below is a polynomial abbreviation:
\begin{gather*}
 J=6j+1,\quad q=\beta J+1,\quad W=12w+4,\quad X=qW,\quad Y=v,\\
 E=XY,\quad a=E+Y,\quad A=a+2,\quad H=4a+3,\quad \Delta=a^2+H,\\
 F=(K_0+W)C-z(q-1),\\
 R=(q(q-F)-C)(q^2-1)+(M_C+qM_F)J,\\
 k=R+1+hE,\quad \eta=c-Yk,\\
 D_{\rm err}=(A_cF-q)^2+(B_cC-q)^2,\\
 B_0=k^2(q^2-D_{\rm err})-4q^2\eta^2,\quad 
 \mathcal B=Y^2B_0-k^2q^8,\\
 \ell=XY^2,\quad L=\ell k,\quad D_\tau=1+L(L+k),\\
 D=3\mathcal B-4,\quad D_g=1+4DD_\tau s^2,\quad P=D_\tau D_g,\\
 D_m=1+\Delta c^2,\quad U_b=rc^2,\quad U=AU_b,\\
 D_f=1+\Delta U^2D_m,\quad Q=\Delta^2U^2D_m,\quad \delta=Q-1,\\
 V=D_f(c-R+4UD_mT)-c,\quad C_y=QV^2-1,\quad D_a=\delta C_y,\\
 K=X(3c-2X)+2,\quad M=2q^3,\quad D_c=HM,\quad S=srHM^5,\\
 z_c=D_cT+MK+H(Y-q^3),\\
 \mathcal A=(z_c^2+3D_c^2)(z_c^3-19D_c^3),\\
 \lambda=(2\Delta+1)U_bD_m,\quad Z=S(n+\lambda V)-\mathcal A.
\end{gather*}
The final equation is the ordinary integer polynomial
\begin{equation}\label{eq:integer352}
 \boxed{\bigl(\delta Z^2+\delta S^2P-4S^2D_mD_fC_y\bigr)^2
             -4\delta^2S^2PZ^2=0.}
\end{equation}
Denote the left-hand side of~\eqref{eq:integer352} by $\Phi_{352}$.
Its literal evaluator is
\path{integer/free_main_saturated11_sharp.json}.
It has no division, radicals, variable exponents, or hidden witnesses.
The displayed radicals in the proof are only an analysis of this polynomial.

\subsection{Why one main square root can replace a full factor}
Unconditionally, $q$ is even, $q\equiv2\pmod3$, $|q|\ge\beta-1$,
$8\mid X$, $X\equiv2\pmod3$, and $|X|\ge4|q|$.
The integer $H$ is odd and nonzero. If $Y=0$, then $\Delta=3$;
otherwise $|a|\ge|X|-1$, so in either case $\Delta=A^2-1>1$.
Thus $D_m,D_f>0$. Also $4\mid Q$, $D_f\equiv1\pmod4$,
$\delta$ is odd and nonzero, and $D_\tau,P$ are odd, with $D_\tau\ge1$.
These parity assertions do not assume $Y$ even: $\Delta A$ is always even.

The intersective quintic $\pi(t)=(t^2+3)(t^3-19)$ has no rational root,
so $\mathcal A=D_c^5\pi(z_c/D_c)\ne0$. If $S=0$, the output is
$\delta^2\mathcal A^4\ne0$. Every zero therefore has $r,s\ne0$.
If $c=0$, then $D_a=1$. If $c\ne0$, then $|U|\ge2c^2$ and
$D_f>|c|+1$; the integer-distance argument in
$V=D_fb-c$ gives $V\ne0$, $Q>1$, and $D_a>0$.

Equation~\eqref{eq:integer352} is $\delta^2$ times the product of
\[
 Z\pm2S\sqrt{D_mD_fC_y/\delta}\pm S\sqrt P.
\]
The first radicand equals $D_mD_fD_a/\delta^2>0$.
Thus $P$ cannot be negative at a zero: it would give the sole nonzero
imaginary term in a vanishing factor. Hence $D_g>0$ and $\mathcal B\ge2$.
The stronger budget yields
\[
 k\ne0,\quad B_0>0,\quad |Y|>|q|^3,\quad
 D_{\rm err}<q^2,\quad |c/k-Y|<\tfrac12.
\]
It supplies the authentic common-sign packing bounds and excludes $c=0$.

Multiplying a normalized factor by $\delta$ gives a rational signed sum
of $2\sqrt{D_mD_fD_a}$ and $\delta\sqrt P$.
Since $P$ and $\delta$ are odd, the two terms have different 2-adic
orders if they belong to the same nontrivial squareclass; they cannot
cancel. Distinct squareclasses are linearly independent. Thus both
$D_mD_fD_a$ and $P$ are squares.

Crucially, $D_m,D_f,D_a$ are pairwise coprime. Modulo $D_m$,
$D_f=D_a=1$. Modulo $D_f$,
\[
 Q=-\Delta,\quad\delta=-A^2,\quad V=-c,\quad C_y=-D_m,
 \quad D_a=A^2D_m.
\]
Because $A\mid U$, $D_f$ is coprime to $A$ as well as $D_m$.
Consequently
\[
 d^2=D_m,\quad f^2=D_f,\quad e^2=D_a,\quad
 \tau^2=D_\tau,\quad g^2=D_g
\]
for positive integers; the last split uses $\gcd(D_\tau,D_g)=1$.
The new integer coordinates used only in the proof are
\[
 I=Ud,\qquad T_{\rm old}=dT.
\]
They satisfy $c^2\mid I$ and restore exactly
\[
 f^2-\Delta I^2=1,\quad Q=(\Delta I)^2,\quad
 V=f^2(c-R+4IT_{\rm old})-c.
\]
The rational $y=e/\delta$ is integral by adjacent-Pell denominator
descent: its norm-preserving transformation decreases $|V|$ while
preserving a reduced denominator, ending at $|V|=|y|=1$.
The factor equation now gives $S\mid\mathcal A$.

\subsection{Late scale recovery without early coprimality}
The divisibility $M^5\mid\mathcal A$ forces $M\mid z_c$ by the
unique lowest valuation of $z_c^5$ when $v_p(z_c)<v_p(M)$.
Hence $M\mid H(Y-q^3)$. Its 2-primary part proves $Y$ even, but
we do not yet cancel $H$ modulo $q$.
Also $H\mid\mathcal A$ gives $H\mid(MK)^5$. With
\[
 K_q=X(X+2ac)-1,\qquad K=HXc-2K_q,
\]
oddness of $H$ gives $H\mid(q^3K_q)^5$.

The retained signed pre-projection rank argument now has every required
hypothesis: the cubic-scale magnitude, centered strip, common signs,
first and main norms, $c^2\mid I$, and the two exact auxiliary residues
$V\equiv-c\pmod f$, $V\equiv-R\pmod c$.
Its nonodd strong-rank step gives $c\mid m$ for the strong Pell index;
it does not assume that $I/c^2$ is odd. It yields
\[
 R=2t+1,\quad t\ge36,\quad
 c=\operatorname{sign}(c)\psi_A(R),\quad
 |H|>|X|(|X|-3)^t.
\]
Here $t\ge36$ follows already from
$R>|q|(|q|^2-1)\ge4080$.

Put $x=|X|$, $\sigma=\operatorname{sign}(c)$ and
$\Pi=(\sigma X-2^R)(2^R\sigma X+1)$.
The Pell projector congruence implies $H\mid q^{15}\Pi^5$.
But
\[
 |q|^{15}|\Pi|^5<2^{20t-10}x^{25},\qquad
 |H|>x^{t+1}/2^t.
\]
The unconditional $x\ge4(\beta-1)\ge2^{64}$ and $t\ge36$ give
$64(t-24)-(21t-10)=43t-1526\ge22>0$.
Thus the divisible integer is smaller than $|H|$ and must be zero.
Its odd second factor is nonzero, so $\sigma X=2^R$.
Now $q\mid X$ makes $|q|$ a power of two. Since $H$ is odd,
$\gcd(H,M)=1$ finally holds, yielding
\[
 Y=(2v_{\rm old}+1)q^3.
\]
The source's orientation, centered half-binomial interval, and fresh
odd-kernel restoration apply. The full proof explicitly constructs the
fresh strong/auxiliary witnesses; it does not relabel incoming $I/c^2$
as odd. An actual old signed Q41 zero is then passed to the retained
native/compiler/refutation theorem.

\subsection{Completeness with one fixed strong pair}
Start with a canonical native witness. It has $D_m=d^2$, $D_\tau=\tau^2$,
$B_0>q^2k^2/4$ and $Y>X=2^R>4q^3$, so the stronger budget exceeds $64$.
Since $D D_\tau$ is positive and $2$ modulo $3$, ordinary Pell existence and
squaring a unit give the required $s,g>0$.
Choose one strong pair with
\[
 4sAdc^2\mid I,\qquad f^2-\Delta I^2=1,
 \qquad r=I/(Adc^2).
\]
All primes of $dsr$ divide this fixed $I$.
The retained supported-modulus interpolation theorem realizes every
$T_{\rm old}$ class modulo $dsr$ without changing this strong pair.
Write $\kappa=K/H$. Choose an intersective root $a_0$ of
$\pi(T+\kappa+v_{\rm old})$ modulo $sr$, and prescribe
\[
 T_{\rm old}=da_0\pmod{dsr}.
\]
Then $T=T_{\rm old}/d$ is integral and
$\mathcal A/S=H^4\pi(T+\kappa+v_{\rm old})/(sr)$ is integral.
Finally set
\[
 n=\mathcal A/S+2dfy+\tau g-\lambda V.
\]
A conjugate factor vanishes. The order is outer data, guard, fixed strong
pair, intersective residue, auxiliary interpolation, and $n$.
This proves both directions of the ordinary zero-existence reduction.

\subsection{Exact degree, evaluation, and verification}
The important degrees are
\[
 \deg(S,D_m,D_f,Q,V,C_y,\lambda,\mathcal A)
       =(20,8,26,32,41,114,17,35).
\]
Writing $W_0=Sn-\mathcal A$, the sharp identity
\[
 \lambda^2-4D_mD_f=D_m^2U_b^2-4D_m
\]
shows that the inner bracket has unique leading term
$S^2QD_m^2U_b^2V^2$, of degree $176$. The final correction has degree $316$,
so the exact output degree is $352$. An attained monomial is
\[
 [j^{124}w^{64}v^{64}s^4r^{24}c^{68}T^4]\Phi_{352}
      =2^{288}3^{188}\beta^{124}>0.
\]
The graph shares the exact expressions
\[
 \mathcal B=(Yk)^2(q^2-D_{\rm err})
          -q^2\bigl(4(Y\eta)^2+(kq^3)^2\bigr),\quad L=E(Yk),
 \qquad P=D_\tau+(12\mathcal B-16)(D_\tau s)^2.
\]
Every scalar multiplication and square is counted.
The direct and simple centers give $376/127$ and $364/131$.
A weaker-scale companion with $Y=qv$, the old budget and $M=2q^2$
gives $428/123$; its detailed early-CRT proof is in the package.

The local checker, without importing the builder, verifies every named
register over free polynomial rings, all four conjugates, coprimality
remainders, rescaling, exact cancellation cuts and attained gate leaders,
topology/liveness, all six numeral payloads and their complete finite
reconstruction, and rejects deliberate corruptions. It passes normally
and with Python optimization. The retained source files and notices are
hashed; their unbounded theorems are an explicit dependency, not something
the finite checks claim to certify.
